\documentclass[10pt,a4paper]{amsart}
\usepackage[margin=19mm]{geometry}
\usepackage{amsmath,amssymb,mathtools}
\usepackage[scr=rsfs,cal=euler]{mathalfa}
\usepackage{xfrac}
\usepackage[unicode,hidelinks,bookmarks=false]{hyperref}
\newcommand{\ie}{i.e.\ }
\newcommand{\cf}{cf.\ }
\newcommand{\resp}{respectively\ }
\newcommand{\Subsection}{\subsection}
\providecommand{\nobreakdash}{\nobreak\hbox{-}}
\setlength{\emergencystretch}{2em}
\multlinegap0pt
\usepackage{bm}
\usepackage{bbm}
\usepackage{dsfont}
\usepackage{xspace}
\usepackage{cancel}
\usepackage{mathtools}
\usepackage[shortlabels]{enumitem}
\usepackage{amsthm}
\makeatletter
\@ifundefined{theorem}{\newtheorem{theorem}{Theorem}}{}
\@ifundefined{lemma}{\newtheorem{lemma}[theorem]{Lemma}}{}
\@ifundefined{proposition}{\newtheorem{proposition}[theorem]{Proposition}}{}
\@ifundefined{remark}{\newtheorem{remark}[theorem]{Remark}}{}
\makeatother
\newcommand{\WAM}{\mathrm{WAM}}  % weak amenability constant (quantitative)
\newcommand{\Zc}{\mathrm{Z}} % centre of a bimodule (avoid \mathbb{Z} clash if you use it)
\newcommand{\ad}{\mathrm{ad}}
\newcommand{\cF}{\mathcal{F}}
\renewcommand{\ge}{\geqslant}
\renewcommand{\le}{\leqslant}
\title[Amenability constants for unconditional sums of Banach algeb]{Amenability constants for unconditional sums of Banach algebras}
\author{Tomasz Kania, Jerzy K\k{a}kol}
\date{Source: arXiv:2601.06680v2 (CC BY 4.0)}
\begin{document}
\maketitle

\noindent\emph{Source and licence.} Amenability constants for unconditional sums of Banach algebras. Version arXiv:2601.06680v2 (CC BY 4.0), distributed under the Creative Commons Attribution 4.0 International licence, as recorded in the arXiv licence metadata for this exact version and shown on the abstract page https://arxiv.org/abs/2601.06680v2. Full rights notice: licenses/arxiv-2601.06680-CC-BY-4.0.txt.\par\smallskip

\noindent\textit{Editorial scope.} Counted unit: the Theorem whose source label is \texttt{thm:WA-counterpart} in the article (source0.tex lines 1471--1483, source label \texttt{thm:WA-counterpart}), delivered together with the article's own complete original proof of it (source0.tex lines 1485--1549, one \texttt{proof} environment of 65 lines). No mathematics is written, repaired or re-derived: statement and proof are the article's own text line for line. Required context retained uncounted: lem:dense-finite-support (lines 524--531); rem:WA-commutative (lines 1356--1359); lem:WA-essential (lines 1375--1377); prop:WA-to-summands (lines 1408--1415); lem:off-diagonal (lines 1441--1450). Every definition, display and label of those lines is retained unchanged. Omitted from this package and not redistributed: 1--523, 532--1355, 1360--1374, 1378--1407, 1416--1440, 1451--1470, 1484--1484, 1550--1777. Nothing inside a retained excerpt is omitted or altered: every retained body is delivered exactly as the article's lines are written, so no omission receipt of any kind accompanies this package. Citations: the retained excerpts carry no citation command at all, so no bibliography entry of the article is transcribed and no reference is redistributed. Reference adaptation: none was needed and none was made; every reference inside the delivered unit resolves inside it, so no unresolved reference or citation is produced. Printed slips kept literal: no printed word, symbol, hypothesis or number of the counted statement or of its proof was altered. Layout and class: the article's own class is \texttt{amsart} with packages including geometry, fontenc, inputenc, amsmath, amssymb, amsthm, mathtools, enumitem, url, hyperref; the wrapper is the lane's pinned \texttt{amsart} editorial wrapper at 10pt on A4, which restates the theorem-like environments the retained text needs and adds the article's own macro definitions that the retained text needs (\texttt{\textbackslash WAM}, \texttt{\textbackslash Zc}, \texttt{\textbackslash ad}, \texttt{\textbackslash cF}, \texttt{\textbackslash ge}, \texttt{\textbackslash le}), with extra wrapper packages (bm, bbm, dsfont, xspace, cancel, mathtools, enumitem). The delivered numbering is the unit's own. Assets: no retained span contains a figure environment, a graphics-inclusion command, a PDF-page inclusion, a table or a TikZ picture; no image file is copied into this package, the package declares no asset dependency and no image file is redistributed with it. Licence: version arXiv:2601.06680v2 of this article is distributed under the Creative Commons Attribution 4.0 International licence, as recorded in the arXiv licence metadata for this exact version and shown on the abstract page https://arxiv.org/abs/2601.06680v2. A full rights notice is delivered at licenses/arxiv-2601.06680-CC-BY-4.0.txt. Characters: every retained excerpt is pure ASCII, so the delivered engine, XeLaTeX with its default Latin Modern text font, reports no missing character.
\begin{lemma}[Finite-support subalgebras are dense]\label{lem:dense-finite-support}
Let $E$ be a Banach sequence lattice on~$I$ and let $(X_i)_{i\in I}$ be Banach spaces.
For $F\in\cF(I)$ let
\[
X_F:=\{x\in (\bigoplus_{i\in I}X_i)_E : x_i=0 \text{ for } i\notin F\}.
\]
Then $\bigcup_{F\in\cF(I)}X_F$ is dense in $(\bigoplus_{i\in I}X_i)_E$.
\end{lemma}

\begin{remark}\label{rem:WA-commutative}
If $A$ is commutative, then $a\cdot\phi=\phi\cdot a$ for all $a\in A$ and $\phi\in A^*$, hence $\ad_\phi=0$ for all $\phi$.
Therefore, for commutative $A$, weak amenability is equivalent to the vanishing of all bounded derivations $A\to A^*$.
\end{remark}

\begin{lemma}\label{lem:WA-essential}
If $A$ is weakly amenable, then $\overline{A^2}=A$, where $A^2:=\mathrm{span}\{ab:\ a,b\in A\}$.
\end{lemma}

\begin{proposition}\label{prop:WA-to-summands}
If $A=(\oplus_{i\in I}A_i)_E$ is weakly amenable, then each $A_i$ is weakly amenable.

More quantitatively, if $\WAM(A)<\infty$, then for each $i\in I$,
\[
\WAM(A_i)\le \|\iota_i\|\,\WAM(A)=\|\delta_i\|_E\,\WAM(A).
\]
\end{proposition}

\begin{lemma}[Off-diagonal matrix coefficients]\label{lem:off-diagonal}
Let $A=(\oplus_{i\in I}A_i)_E$ and let $D\colon A\to A^*$ be a bounded derivation.
For $j,k\in I$ define
\[
D_{jk}:=\iota_j^*\circ D\circ\iota_k\colon A_k\longrightarrow A_j^*.
\]
If $j\neq k$ and $\overline{A_j^2}=A_j$, then $D_{jk}=0$.  Consequently, if every $A_j$ is
weakly amenable, then for each $k\in I$ and $a\in A_k$ the functional
$D(\iota_k(a))$ annihilates every coordinate ideal $\iota_j(A_j)$ with $j\neq k$.
\end{lemma}

\begin{theorem}\label{thm:WA-counterpart}
Let $I$ be a non-empty set, let $E$ be a Banach sequence lattice on $I$, and let $(A_i)_{i\in I}$ be Banach algebras.
Set $A:=(\oplus_{i\in I}A_i)_E$.

\begin{enumerate}[label={\rm(\arabic*)},leftmargin=3.2em]
\item If $A$ is weakly amenable, then each $A_i$ is weakly amenable.
\item If every $A_i$ is commutative and weakly amenable, then $A$ is weakly amenable.
(In fact every derivation $A\to A^*$ is zero.)
\item Assume $I$ is infinite and $E=\ell_p(I)$ for some $1<p<\infty$.
Let $B$ be a weakly amenable \emph{non-commutative} Banach algebra, and form the constant family $A_i=B$.
Then the algebra $(\oplus_{i\in I}B)_{\ell_p(I)}$ is \emph{not} weakly amenable.
\end{enumerate}
\end{theorem}

\begin{proof}
\emph{(1)} This is Proposition~\ref{prop:WA-to-summands}.

\smallskip

\emph{(2)} Let $D\colon A\to A^*$ be a bounded derivation.
Since $A$ is commutative, inner derivations vanish (Remark~\ref{rem:WA-commutative}), so it suffices to prove $D=0$.

By Lemma~\ref{lem:WA-essential}, each summand satisfies $\overline{A_i^2}=A_i$.
Hence Lemma~\ref{lem:off-diagonal} shows that, for $j\in I$ and $a\in A_j$, the
functional $D(\iota_j(a))$ annihilates every coordinate ideal $\iota_i(A_i)$ with
$i\neq j$.

It remains to look at the diagonal coefficient.  Define
\[
d_j:=\iota_j^*\circ D\circ \iota_j\colon A_j\to A_j^*.
\]
Then $d_j$ is a derivation.  Since $A_j$ is commutative and weakly amenable, $d_j=0$.
Thus $D(\iota_j(a))$ also annihilates $\iota_j(A_j)$.  Consequently
$D(\iota_j(a))$ annihilates the algebraic span of all coordinate ideals, which is dense in
$A$ by Lemma~\ref{lem:dense-finite-support}; by continuity of the functional
$D(\iota_j(a))$, it is zero.  We have shown that $D$ vanishes on every coordinate ideal,
hence on the dense subalgebra of finitely supported elements.  Since $D$ is continuous,
$D=0$ on $A$.

\smallskip
\emph{(3)} Let $1<p<\infty$ and put $q=p/(p-1)$.
Choose $\psi\in B^*$ with $\psi\notin \Zc(B^*)$; equivalently $\ad_\psi\neq 0$.
Set $d:=\mathrm{dist}(\psi,\Zc(B^*))>0$.

Choose a countably infinite subset $\{i_n:n\in\mathbb N\}\subset I$.
Define a scalar weight $w=(w_i)_{i\in I}$ by
\[
w_{i_n}:=
\begin{cases}
1, & 1<p\le 2,\\[2pt]
n^{-1/q}, & p>2,
\end{cases}
\qquad\text{and}\qquad
w_i:=0\ \ (i\notin \{i_n:n\in\mathbb N\}).
\]
Then $w\notin \ell_q(I)$, but multiplication by $w$ defines a bounded operator
$M_w\colon \ell_p(I)\to \ell_q(I)$:
for $1<p\le 2$ this is just the continuous embedding $\ell_p\hookrightarrow \ell_q$;
for $p>2$ it follows from H\"older since $w\in \ell_r(I)$ with $r=p/(p-2)$ and $1/q=1/p+1/r$.

Now define $D\colon \ell_p(I,B)\to \ell_q(I,B^*)\cong (\ell_p(I,B))^*$ by
\[
D\bigl((b_i)_{i\in I}\bigr):=\bigl(w_i\,\ad_\psi(b_i)\bigr)_{i\in I}.
\]
Since $\|\ad_\psi(b_i)\|\le 2\|\psi\|\,\|b_i\|$, boundedness of $M_w$ yields that $D$ is a bounded derivation.

Suppose, towards a contradiction, that $D$ is inner: $D=\ad_\Phi$ for some $\Phi\in \ell_q(I,B^*)$.
Restricting to the $i$th coordinate ideal (via the canonical embedding/projection) gives
\[
\ad_{\Phi_i}=w_i\,\ad_\psi\qquad (i\in I),
\]
hence $\Phi_i-w_i\psi\in \Zc(B^*)$ for each $i$.
Therefore
\[
\|\Phi_i\|\ge \mathrm{dist}(w_i\psi,\Zc(B^*))=|w_i|\,\mathrm{dist}(\psi,\Zc(B^*))=d\,|w_i|.
\]
It follows that $(\|\Phi_i\|)_{i\in I}$ dominates $d(|w_i|)_{i\in I}$ and hence cannot lie in $\ell_q(I)$ because $w\notin\ell_q(I)$.
This contradicts $\Phi\in\ell_q(I,B^*)$. Thus $D$ is not inner and $\ell_p(I,B)$ is not weakly amenable.
\end{proof}

\end{document}
